\documentclass[11pt]{article}
% The preamble every project's paper shares: % The preamble every project's paper shares: % The preamble every project's paper shares: \input{../../preamble} from
% projects/<project>/paper/. Boards come from chessboard and xskak, plots from
% pgfplots; every number a paper prints comes from its generated.tex.
\usepackage[LSBC4,T1]{fontenc}
\usepackage[utf8]{inputenc}
\usepackage{lmodern}
\usepackage[margin=1in]{geometry}
\usepackage{microtype}
\usepackage{amsmath}
\usepackage{booktabs}
\usepackage{tabularx}
\usepackage{array}
\usepackage{xcolor}
\usepackage[obeyspaces,spaces]{url}
\usepackage{xskak}
\usepackage{pgfplots}
\pgfplotsset{compat=1.18}
\usepackage{caption}
\usepackage{subcaption}
\usepackage{enumitem}
\usepackage[section]{placeins}
\usepackage[hidelinks]{hyperref}

\setchessboard{boardfontencoding=LSBC4, boardfontfamily=skaknew, blackfieldcolor=black!22,
  whitefieldcolor=white, setfontcolors, showmover=false, label=true, labelfontsize=6pt,
  boardfontsize=14pt}

% A move or line in UCI, and a position in FEN, both breakable at a slash.
\newcommand{\uci}[1]{\texttt{#1}}
\newcommand{\fen}[1]{{\footnotesize\path{#1}}}
\newcolumntype{L}{>{\raggedright\arraybackslash}X}

\setlength{\parskip}{0.4em}
\setlength{\parindent}{0pt}
 from
% projects/<project>/paper/. Boards come from chessboard and xskak, plots from
% pgfplots; every number a paper prints comes from its generated.tex.
\usepackage[LSBC4,T1]{fontenc}
\usepackage[utf8]{inputenc}
\usepackage{lmodern}
\usepackage[margin=1in]{geometry}
\usepackage{microtype}
\usepackage{amsmath}
\usepackage{booktabs}
\usepackage{tabularx}
\usepackage{array}
\usepackage{xcolor}
\usepackage[obeyspaces,spaces]{url}
\usepackage{xskak}
\usepackage{pgfplots}
\pgfplotsset{compat=1.18}
\usepackage{caption}
\usepackage{subcaption}
\usepackage{enumitem}
\usepackage[section]{placeins}
\usepackage[hidelinks]{hyperref}

\setchessboard{boardfontencoding=LSBC4, boardfontfamily=skaknew, blackfieldcolor=black!22,
  whitefieldcolor=white, setfontcolors, showmover=false, label=true, labelfontsize=6pt,
  boardfontsize=14pt}

% A move or line in UCI, and a position in FEN, both breakable at a slash.
\newcommand{\uci}[1]{\texttt{#1}}
\newcommand{\fen}[1]{{\footnotesize\path{#1}}}
\newcolumntype{L}{>{\raggedright\arraybackslash}X}

\setlength{\parskip}{0.4em}
\setlength{\parindent}{0pt}
 from
% projects/<project>/paper/. Boards come from chessboard and xskak, plots from
% pgfplots; every number a paper prints comes from its generated.tex.
\usepackage[LSBC4,T1]{fontenc}
\usepackage[utf8]{inputenc}
\usepackage{lmodern}
\usepackage[margin=1in]{geometry}
\usepackage{microtype}
\usepackage{amsmath}
\usepackage{booktabs}
\usepackage{tabularx}
\usepackage{array}
\usepackage{xcolor}
\usepackage[obeyspaces,spaces]{url}
\usepackage{xskak}
\usepackage{pgfplots}
\pgfplotsset{compat=1.18}
\usepackage{caption}
\usepackage{subcaption}
\usepackage{enumitem}
\usepackage[section]{placeins}
\usepackage[hidelinks]{hyperref}

\setchessboard{boardfontencoding=LSBC4, boardfontfamily=skaknew, blackfieldcolor=black!22,
  whitefieldcolor=white, setfontcolors, showmover=false, label=true, labelfontsize=6pt,
  boardfontsize=14pt}

% A move or line in UCI, and a position in FEN, both breakable at a slash.
\newcommand{\uci}[1]{\texttt{#1}}
\newcommand{\fen}[1]{{\footnotesize\path{#1}}}
\newcolumntype{L}{>{\raggedright\arraybackslash}X}

\setlength{\parskip}{0.4em}
\setlength{\parindent}{0pt}

% Written by projects/winprob-parity/paper/generate.py; do not edit by hand.

\newcommand{\AppCommit}{\texttt{d8b7712}}
\newcommand{\GridMinPawns}{-120}
\newcommand{\GridMaxPawns}{120}
\newcommand{\GridStepPawns}{0.1}
\newcommand{\MateMax}{30}
\newcommand{\MateStep}{10}
\newcommand{\FigurePawns}{30}
\newcommand{\GapAxisMax}{2.7}
\newcommand{\AppClampPawns}{100}
\newcommand{\AppClamp}{10000}
\newcommand{\TolerancePower}{-9}
\newcommand{\WonWithin}{0.01}
\newcommand{\WonAtPawns}{25.5}
\newcommand{\CrossLoss}{95.09}
\newcommand{\Slope}{0.00368208}
\newcommand{\LabClamp}{1000}
\newcommand{\LabClampPawns}{10}
\newcommand{\MateCp}{10000}
\newcommand{\BlunderThreshold}{20}
\newcommand{\CandidateThreshold}{10}
\newcommand{\SharedPoints}{201}
\newcommand{\SharedMaxGap}{0}
\newcommand{\PastPoints}{1,800}
\newcommand{\BeyondPoints}{400}
\newcommand{\MatePoints}{60}
\newcommand{\AllPoints}{2,461}
\newcommand{\MaxGap}{2.46}
\newcommand{\LabCeiling}{97.54}
\newcommand{\LabFloor}{2.46}
\newcommand{\SharedAgrees}{agree}
\newcommand{\GridPoints}{481}
\newcommand{\GridPairs}{231,361}
\newcommand{\GridDiffer}{872}
\newcommand{\GridBlunderToCandidate}{436}
\newcommand{\GridCandidateToNone}{436}
\newcommand{\GridWidest}{2.46}
\newcommand{\ExampleFrom}{$+$15.00}
\newcommand{\ExampleTo}{$+$3.50}
\newcommand{\ExampleAppBefore}{99.60}
\newcommand{\ExampleLabBefore}{97.54}
\newcommand{\ExampleAfter}{78.39}
\newcommand{\ExampleAppLoss}{21.21}
\newcommand{\ExampleLabLoss}{19.15}
\newcommand{\ParityTable}{%
\begin{tabular}{lrrlrr}
\toprule
range & points & largest gap & at & lab & app \\
\midrule
within $\pm$1000 cp & 201 & 0.0000 & --- & 37.3907 & 37.3907 \\
$\pm$1000 to $\pm$10000 cp & 1,800 & 2.4553 & $+$9980 cp & 97.5447 & 100.0000 \\
beyond $\pm$10000 cp & 400 & 2.4553 & $+$10010 cp & 97.5447 & 100.0000 \\
mate in 1 to 30 & 60 & 2.4553 & \#$+$1 & 97.5447 & 100.0000 \\
\bottomrule
\end{tabular}
}


\title{Win probability in the lab and in the app\\[0.3em]
\large One curve, two clamps, and the move labels they can change}
\author{Chess Coach research lab}
\date{}

\begin{document}
\maketitle

\begin{abstract}
The lab and the coaching app turn an engine evaluation into a win probability with the same
curve, the logistic Lichess fits on its games, with a slope of \Slope{} per centipawn. They
clamp it differently: the lab clamps the evaluation at $\pm\LabClampPawns$ pawns and reads
every mate as that clamp, while the app clamps at $\pm\AppClampPawns$ pawns and reads mate in
$n$ as $\pm(\AppClamp - \MateStep n)$ centipawns. We evaluated both at \AllPoints{} evaluations, the app's
through its own code at commit \AppCommit. Within $\pm\LabClampPawns$ pawns the two agree at
every point. Past it the lab holds every evaluation, mates included, at \LabCeiling\% for the
side ahead, while the app goes on to 100\%: a gap of up to \MaxGap{} points. The lab's move
labels (a loss of more than \CandidateThreshold{} or \BlunderThreshold{} points of the mover's
win probability) can therefore differ from the same labels read on the app's curve, but only
for a move that starts more than \LabClampPawns{} pawns ahead or ends more than
\LabClampPawns{} pawns behind, and only by one step, the lab's being the milder.
\end{abstract}

\section{Why the two must agree}

The lab trains its models on labels read from win probability. Its data pipeline marks a move
as a blunder when it costs the mover more than \BlunderThreshold{} points of win probability
(the blunder hazard model's target, and the early traps the repertoire optimizer counts), and
the graded-opponent scorer keeps a move as a candidate when it costs more than
\CandidateThreshold. A model that wins in the lab is ported to the app, and the app computes its
own verdicts from its own win probability. If the two curves differ, a label the lab trained on
means something else once the model runs in the app.

The lab never re-implements an app instrument to compare against it. This check calls the
app's own function (\texttt{winPercent} in \texttt{src/core/score}) at the commit the lab pins,
through the lab's bridge, and compares it with the lab's (\texttt{win\_percent} in
\texttt{chess\_coach\_models.winprob}).

\section{One curve, two clamps}

Both sides compute White's win probability from an evaluation of $x$ centipawns from White's
side as
\[
  W(x) \;=\; 50 + 50\left(\frac{2}{1+e^{-kx}} - 1\right),\qquad k = \Slope,
\]
and differ only in what they feed it.

\begin{itemize}[nosep]
\item \textbf{The lab} clamps $x$ to $\pm\LabClamp$ before the curve, as Lichess does. A mate
from a game's annotations, or from the engine, becomes $\pm\MateCp$ (or within a few
centipawns of it) and is then clamped like any other evaluation, so every mate reads
\LabCeiling\% for the side delivering it and \LabFloor\% for the side receiving it.
\item \textbf{The app} clamps $x$ to $\pm\AppClamp$ and reads mate in $n$ as $\pm(\AppClamp - \MateStep n)$,
so a mate reads as won, a little less won the further off it is. Its code notes that the wider clamp
is the one every number it has stored was computed with. A finished game or a tablebase result
reads exactly 100, 50 or 0; the lab has no such reading, and this check leaves it out.
\end{itemize}

\section{Method}

Both functions were evaluated at every evaluation from $\GridMinPawns$ to $+\GridMaxPawns$ pawns
in steps of \GridStepPawns{} pawns and at mate in 1 to \MateMax{} for either side: \AllPoints{} points, \SharedPoints{} of
them within the lab's clamp. The lab's side runs in Python, the app's in Node under the app's
own toolchain at \AppCommit{} (\texttt{parity.py} sends the evaluations to \texttt{app.mts} on
standard input and reads White's win probability back). The check fails if any point within
$\pm\LabClampPawns$ pawns differs by more than $10^{\TolerancePower}$ points, the room left for the two
languages' exponentials to differ in the last bit. Past the clamp it reports, for each range,
the largest gap and where it falls.

\section{Results}

Within $\pm\LabClampPawns$ pawns the two curves \SharedAgrees{} at all \SharedPoints{} points
(the largest difference rounds to \SharedMaxGap{} at six decimals). Past the clamp they part as
designed (Table~\ref{tab:parity}, Figure~\ref{fig:curves}): the lab's curve is flat at
\LabCeiling\% while the app's keeps rising, and from \WonAtPawns{} pawns on the app reads the side
ahead as won to within \WonWithin{} points. Every mate shows the full gap, \MaxGap{} points, because
the lab reads it at the clamp and the app at 100\% less a sliver.

\begin{table}[!htbp]
\centering\small
\ParityTable
\caption{White's win probability in the lab and in the app, in percentage points: the largest gap
in each range, the first evaluation where it occurs, and the two readings there. Within the lab's
clamp the gap is zero at every point.}
\label{tab:parity}
\end{table}

\begin{figure}[!htbp]
\centering
\begin{subfigure}{0.49\textwidth}\centering
\begin{tikzpicture}
\begin{axis}[scale only axis, width=0.78\linewidth, height=4.4cm,
  xlabel={evaluation, pawns (White's side)}, ylabel={White's win \%}, xmin=0, xmax=\FigurePawns, ymin=50, ymax=101, grid=major,
  legend pos=south east, legend cell align=left]
\addplot[thick, black] table[col sep=comma, x expr=\thisrow{cp}/100, y=app,
  restrict x to domain=0:\FigurePawns] {../results/parity_winprob_curve.csv};
\addlegendentry{app}
\addplot[thick, red!70!black, dashed] table[col sep=comma, x expr=\thisrow{cp}/100, y=lab,
  restrict x to domain=0:\FigurePawns] {../results/parity_winprob_curve.csv};
\addlegendentry{lab}
\draw[gray, densely dotted] (axis cs:\LabClampPawns,50) -- (axis cs:\LabClampPawns,101);
\end{axis}
\end{tikzpicture}
\caption{The two curves. The dotted line is the lab's clamp.}
\end{subfigure}\hfill
\begin{subfigure}{0.49\textwidth}\centering
\begin{tikzpicture}
\begin{axis}[scale only axis, width=0.78\linewidth, height=4.4cm,
  xlabel={evaluation, pawns (White's side)}, ylabel={app $-$ lab, points}, xmin=0, xmax=\FigurePawns, ymin=0, ymax=\GapAxisMax, grid=major]
\addplot[thick, black] table[col sep=comma, x expr=\thisrow{cp}/100,
  y expr=\thisrow{app}-\thisrow{lab}, restrict x to domain=0:\FigurePawns]
  {../results/parity_winprob_curve.csv};
\end{axis}
\end{tikzpicture}
\caption{The gap. It is zero up to the clamp and approaches \MaxGap{} points past it.}
\end{subfigure}
\caption{White's win probability for White's evaluations from level to \FigurePawns{} pawns ahead, every half
pawn, both columns written by \texttt{parity.py}. Black's side is the mirror image.}
\label{fig:curves}
\end{figure}

\section{Which labels can change}

A move's label depends on its loss, the mover's win probability before the move less the one
after. Read from the mover's side, the two curves give the same loss when both evaluations lie
within the clamp, and the same label when both lie past it on the same side (a loss of at most
\MaxGap{} points either way, too small for any label) or on opposite sides (a change of more than
\CrossLoss{} points on either curve: a blunder on both, or a gain on both). A label can differ only when one evaluation lies
past the clamp and the other within it, and of those only a move that starts more than
\LabClampPawns{} pawns ahead, or ends more than \LabClampPawns{} pawns behind, has a positive loss
to label. There the lab's loss is smaller than the app's by the gap between the two readings of
the evaluation past the clamp, at most \MaxGap{} points, so the lab's label is one step milder
whenever the app's loss exceeds a threshold by no more than that gap.

\paragraph{The claim, checked.} Over every move from one evaluation to another on a grid of
\GridPoints{} evaluations (every half pawn from $\GridMinPawns$ to $+\GridMaxPawns$; \GridPairs{}
pairs, both curves as \texttt{parity.py} wrote them), \GridDiffer{} get different labels: \GridBlunderToCandidate{}
blunders on the app's curve are candidates in the lab, and \GridCandidateToNone{} candidates on
the app's curve are unlabelled in the lab. Every one starts more than \LabClampPawns{} pawns ahead
for the mover or ends more than \LabClampPawns{} pawns behind, in every one the lab's label is one
step milder, and no window is wider than \GridWidest{} points. \texttt{generate.py} checks all
three and refuses to write this paper if any fails. The counts say how many grid pairs are
affected, not how many moves in real games are.

\paragraph{An example.} A player who is \ExampleFrom{} pawns up plays a move that leaves them
\ExampleTo{} pawns up, the smallest such drop the app's curve calls a blunder. The app reads the
position before at \ExampleAppBefore\% and after at \ExampleAfter\%, a loss of
\ExampleAppLoss{} points: a blunder. The lab reads the position before at \ExampleLabBefore\% and
after at the same \ExampleAfter\%, a loss of \ExampleLabLoss{} points: a candidate, not a
blunder. In a won position, the lab forgives a little more.

\section{What a port has to decide}

Which clamp a ported model uses is decided in its port, not here. A model whose labels or
features were computed with the lab's clamp either keeps that clamp inside the port and says so,
or is retrained on labels read from the app's curve. This check is how either choice stays
honest: it runs again whenever the lab moves its pin to a new app commit (\texttt{make parity}),
and fails if the curves ever differ within the clamp.

\section{Reproduction}

\texttt{parity.py} writes the report (\texttt{results/parity\_winprob.json} and \texttt{.md}) and
the curve every half pawn (\texttt{results/parity\_winprob\_curve.csv}) from the app at
\AppCommit; \texttt{paper/generate.py} reads those files, the lab's thresholds and the lab's own
curve, and writes every number this paper prints. The app's repository is private, so the check
runs where the bridge can read it, and its results are committed with the commit they measured.

\section{Limitations}

\begin{itemize}[nosep]
\item The grid is uniform in evaluation. How often real moves fall in the window where a label
differs, on Lichess games at the levels the lab trains on, is not measured here.
\item Only evaluations and mates are compared. A finished game or a tablebase result, which the
app reads exactly, has no counterpart in the lab.
\end{itemize}

\end{document}
